% Erratum E10 for inclusion in main.tex.
% Suggested use: % Erratum E10 for inclusion in main.tex.
% Suggested use: % Erratum E10 for inclusion in main.tex.
% Suggested use: % Erratum E10 for inclusion in main.tex.
% Suggested use: \input{erratum_E10_tikk_endpoints}

\subsection*{Erratum E10: endpoint cases in Proposition 3.2}
\label{s:erratum:E10}

The proof of Proposition 3.2 says that the first inequality in
inequality (59) is a consequence of equation (6) with
$\theta=\nu/k$.  This is correct in the interior case $0<\nu<k$.  At the
endpoints $\nu=0$ and $\nu=k$, however, the displayed choice of the
parameter in equation (6) degenerates, as noted in Erratum E1.  The
statement of Proposition 3.2 remains correct, but the endpoint
cases should be justified separately.

Put
\[
        e:=x^\dag-x_0,
        \qquad
        A:=T^*T,
        \qquad
        S_\alpha:=\alpha(A+\alpha I)^{-1} .
\]
By equation (44),
\[
        x^\dag-x_{\alpha,k}=S_\alpha^k e .
\]
The spectral representation used in the proof is
\begin{align}
\label{e:erratum:E10:spectral}
        \alpha^{-2\nu}\norm{S_\alpha^k e}{}^2
        =
        \alpha^{-2\nu}
        \int_{[0,\infty)}
        \left(\frac{\alpha}{\alpha+\lambda}\right)^{2k}
        d\mu_e(\lambda).
\end{align}

First let $\nu=0$.  Then $N_0=1$ and, by the identity $X_{0:k}=X$ from
equation (25),
\[
        \norm{e}{0:k}=\norm{e}{} .
\]
For every $\alpha>0$ the multiplier
$\alpha/(\alpha+\lambda)$ is bounded by $1$, so
\[
        \norm{S_\alpha^k e}{} \leq \norm{e}{} .
\]
On the other hand,
\[
        \left(\frac{\alpha}{\alpha+\lambda}\right)^{2k}
        \longrightarrow 1
        \qquad (\alpha\to\infty)
\]
for every $\lambda\geq0$, and the integrand is bounded by $1$.  Hence
Lebesgue dominated convergence gives
\[
        \lim_{\alpha\to\infty}\norm{S_\alpha^k e}{}^2
        =
        \int_{[0,\infty)} d\mu_e(\lambda)
        =
        \norm{e}{}^2 .
\]
Thus
\[
        \sup_{\alpha>0}\norm{x^\dag-x_{\alpha,k}}{}
        =
        \norm{e}{}
        =
        \norm{e}{0:k},
\]
which proves inequality (59) at $\nu=0$, in fact with equality.

Now let $\nu=k$.  Then $N_1=1$ and, again by the identity $X_{k:k}=X_k$
from equation (25),
\[
        \norm{e}{k:k}=\norm{e}{k} .
\]
From \eqref{e:erratum:E10:spectral},
\begin{align}
\label{e:erratum:E10:nuk}
        \alpha^{-2k}\norm{S_\alpha^k e}{}^2
        =
        \int_{[0,\infty)}
        (\alpha+\lambda)^{-2k}
        d\mu_e(\lambda).
\end{align}
Since $T$ is injective, $E_{\{0\}}=0$, hence $\mu_e(\{0\})=0$.  For
$\lambda>0$ the function $(\alpha+\lambda)^{-2k}$ increases to
$\lambda^{-2k}$ as $\alpha\searrow0$.  Therefore the monotone convergence
theorem, with the value $+\infty$ allowed, gives
\[
        \lim_{\alpha\searrow0}
        \alpha^{-2k}\norm{S_\alpha^k e}{}^2
        =
        \int_{[0,\infty)} \lambda^{-2k}\,d\mu_e(\lambda)
        =
        \norm{e}{k}^2 .
\]
Consequently,
\[
        \sup_{\alpha>0}
        \alpha^{-k}\norm{x^\dag-x_{\alpha,k}}{}
        =
        \lim_{\alpha\searrow0}
        \alpha^{-k}\norm{x^\dag-x_{\alpha,k}}{}
        =
        \norm{e}{k}
        =
        \norm{e}{k:k} .
\]
This proves inequality (59) at $\nu=k$ and also justifies the final
sentence of Proposition 3.2.


\subsection*{Erratum E10: endpoint cases in Proposition 3.2}
\label{s:erratum:E10}

The proof of Proposition 3.2 says that the first inequality in
inequality (59) is a consequence of equation (6) with
$\theta=\nu/k$.  This is correct in the interior case $0<\nu<k$.  At the
endpoints $\nu=0$ and $\nu=k$, however, the displayed choice of the
parameter in equation (6) degenerates, as noted in Erratum E1.  The
statement of Proposition 3.2 remains correct, but the endpoint
cases should be justified separately.

Put
\[
        e:=x^\dag-x_0,
        \qquad
        A:=T^*T,
        \qquad
        S_\alpha:=\alpha(A+\alpha I)^{-1} .
\]
By equation (44),
\[
        x^\dag-x_{\alpha,k}=S_\alpha^k e .
\]
The spectral representation used in the proof is
\begin{align}
\label{e:erratum:E10:spectral}
        \alpha^{-2\nu}\norm{S_\alpha^k e}{}^2
        =
        \alpha^{-2\nu}
        \int_{[0,\infty)}
        \left(\frac{\alpha}{\alpha+\lambda}\right)^{2k}
        d\mu_e(\lambda).
\end{align}

First let $\nu=0$.  Then $N_0=1$ and, by the identity $X_{0:k}=X$ from
equation (25),
\[
        \norm{e}{0:k}=\norm{e}{} .
\]
For every $\alpha>0$ the multiplier
$\alpha/(\alpha+\lambda)$ is bounded by $1$, so
\[
        \norm{S_\alpha^k e}{} \leq \norm{e}{} .
\]
On the other hand,
\[
        \left(\frac{\alpha}{\alpha+\lambda}\right)^{2k}
        \longrightarrow 1
        \qquad (\alpha\to\infty)
\]
for every $\lambda\geq0$, and the integrand is bounded by $1$.  Hence
Lebesgue dominated convergence gives
\[
        \lim_{\alpha\to\infty}\norm{S_\alpha^k e}{}^2
        =
        \int_{[0,\infty)} d\mu_e(\lambda)
        =
        \norm{e}{}^2 .
\]
Thus
\[
        \sup_{\alpha>0}\norm{x^\dag-x_{\alpha,k}}{}
        =
        \norm{e}{}
        =
        \norm{e}{0:k},
\]
which proves inequality (59) at $\nu=0$, in fact with equality.

Now let $\nu=k$.  Then $N_1=1$ and, again by the identity $X_{k:k}=X_k$
from equation (25),
\[
        \norm{e}{k:k}=\norm{e}{k} .
\]
From \eqref{e:erratum:E10:spectral},
\begin{align}
\label{e:erratum:E10:nuk}
        \alpha^{-2k}\norm{S_\alpha^k e}{}^2
        =
        \int_{[0,\infty)}
        (\alpha+\lambda)^{-2k}
        d\mu_e(\lambda).
\end{align}
Since $T$ is injective, $E_{\{0\}}=0$, hence $\mu_e(\{0\})=0$.  For
$\lambda>0$ the function $(\alpha+\lambda)^{-2k}$ increases to
$\lambda^{-2k}$ as $\alpha\searrow0$.  Therefore the monotone convergence
theorem, with the value $+\infty$ allowed, gives
\[
        \lim_{\alpha\searrow0}
        \alpha^{-2k}\norm{S_\alpha^k e}{}^2
        =
        \int_{[0,\infty)} \lambda^{-2k}\,d\mu_e(\lambda)
        =
        \norm{e}{k}^2 .
\]
Consequently,
\[
        \sup_{\alpha>0}
        \alpha^{-k}\norm{x^\dag-x_{\alpha,k}}{}
        =
        \lim_{\alpha\searrow0}
        \alpha^{-k}\norm{x^\dag-x_{\alpha,k}}{}
        =
        \norm{e}{k}
        =
        \norm{e}{k:k} .
\]
This proves inequality (59) at $\nu=k$ and also justifies the final
sentence of Proposition 3.2.


\subsection*{Erratum E10: endpoint cases in Proposition 3.2}
\label{s:erratum:E10}

The proof of Proposition 3.2 says that the first inequality in
inequality (59) is a consequence of equation (6) with
$\theta=\nu/k$.  This is correct in the interior case $0<\nu<k$.  At the
endpoints $\nu=0$ and $\nu=k$, however, the displayed choice of the
parameter in equation (6) degenerates, as noted in Erratum E1.  The
statement of Proposition 3.2 remains correct, but the endpoint
cases should be justified separately.

Put
\[
        e:=x^\dag-x_0,
        \qquad
        A:=T^*T,
        \qquad
        S_\alpha:=\alpha(A+\alpha I)^{-1} .
\]
By equation (44),
\[
        x^\dag-x_{\alpha,k}=S_\alpha^k e .
\]
The spectral representation used in the proof is
\begin{align}
\label{e:erratum:E10:spectral}
        \alpha^{-2\nu}\norm{S_\alpha^k e}{}^2
        =
        \alpha^{-2\nu}
        \int_{[0,\infty)}
        \left(\frac{\alpha}{\alpha+\lambda}\right)^{2k}
        d\mu_e(\lambda).
\end{align}

First let $\nu=0$.  Then $N_0=1$ and, by the identity $X_{0:k}=X$ from
equation (25),
\[
        \norm{e}{0:k}=\norm{e}{} .
\]
For every $\alpha>0$ the multiplier
$\alpha/(\alpha+\lambda)$ is bounded by $1$, so
\[
        \norm{S_\alpha^k e}{} \leq \norm{e}{} .
\]
On the other hand,
\[
        \left(\frac{\alpha}{\alpha+\lambda}\right)^{2k}
        \longrightarrow 1
        \qquad (\alpha\to\infty)
\]
for every $\lambda\geq0$, and the integrand is bounded by $1$.  Hence
Lebesgue dominated convergence gives
\[
        \lim_{\alpha\to\infty}\norm{S_\alpha^k e}{}^2
        =
        \int_{[0,\infty)} d\mu_e(\lambda)
        =
        \norm{e}{}^2 .
\]
Thus
\[
        \sup_{\alpha>0}\norm{x^\dag-x_{\alpha,k}}{}
        =
        \norm{e}{}
        =
        \norm{e}{0:k},
\]
which proves inequality (59) at $\nu=0$, in fact with equality.

Now let $\nu=k$.  Then $N_1=1$ and, again by the identity $X_{k:k}=X_k$
from equation (25),
\[
        \norm{e}{k:k}=\norm{e}{k} .
\]
From \eqref{e:erratum:E10:spectral},
\begin{align}
\label{e:erratum:E10:nuk}
        \alpha^{-2k}\norm{S_\alpha^k e}{}^2
        =
        \int_{[0,\infty)}
        (\alpha+\lambda)^{-2k}
        d\mu_e(\lambda).
\end{align}
Since $T$ is injective, $E_{\{0\}}=0$, hence $\mu_e(\{0\})=0$.  For
$\lambda>0$ the function $(\alpha+\lambda)^{-2k}$ increases to
$\lambda^{-2k}$ as $\alpha\searrow0$.  Therefore the monotone convergence
theorem, with the value $+\infty$ allowed, gives
\[
        \lim_{\alpha\searrow0}
        \alpha^{-2k}\norm{S_\alpha^k e}{}^2
        =
        \int_{[0,\infty)} \lambda^{-2k}\,d\mu_e(\lambda)
        =
        \norm{e}{k}^2 .
\]
Consequently,
\[
        \sup_{\alpha>0}
        \alpha^{-k}\norm{x^\dag-x_{\alpha,k}}{}
        =
        \lim_{\alpha\searrow0}
        \alpha^{-k}\norm{x^\dag-x_{\alpha,k}}{}
        =
        \norm{e}{k}
        =
        \norm{e}{k:k} .
\]
This proves inequality (59) at $\nu=k$ and also justifies the final
sentence of Proposition 3.2.


\subsection*{Erratum E10: endpoint cases in Proposition 3.2}
\label{s:erratum:E10}

The proof of Proposition 3.2 says that the first inequality in
inequality (59) is a consequence of equation (6) with
$\theta=\nu/k$.  This is correct in the interior case $0<\nu<k$.  At the
endpoints $\nu=0$ and $\nu=k$, however, the displayed choice of the
parameter in equation (6) degenerates, as noted in Erratum E1.  The
statement of Proposition 3.2 remains correct, but the endpoint
cases should be justified separately.

Put
\[
        e:=x^\dag-x_0,
        \qquad
        A:=T^*T,
        \qquad
        S_\alpha:=\alpha(A+\alpha I)^{-1} .
\]
By equation (44),
\[
        x^\dag-x_{\alpha,k}=S_\alpha^k e .
\]
The spectral representation used in the proof is
\begin{align}
\label{e:erratum:E10:spectral}
        \alpha^{-2\nu}\norm{S_\alpha^k e}{}^2
        =
        \alpha^{-2\nu}
        \int_{[0,\infty)}
        \left(\frac{\alpha}{\alpha+\lambda}\right)^{2k}
        d\mu_e(\lambda).
\end{align}

First let $\nu=0$.  Then $N_0=1$ and, by the identity $X_{0:k}=X$ from
equation (25),
\[
        \norm{e}{0:k}=\norm{e}{} .
\]
For every $\alpha>0$ the multiplier
$\alpha/(\alpha+\lambda)$ is bounded by $1$, so
\[
        \norm{S_\alpha^k e}{} \leq \norm{e}{} .
\]
On the other hand,
\[
        \left(\frac{\alpha}{\alpha+\lambda}\right)^{2k}
        \longrightarrow 1
        \qquad (\alpha\to\infty)
\]
for every $\lambda\geq0$, and the integrand is bounded by $1$.  Hence
Lebesgue dominated convergence gives
\[
        \lim_{\alpha\to\infty}\norm{S_\alpha^k e}{}^2
        =
        \int_{[0,\infty)} d\mu_e(\lambda)
        =
        \norm{e}{}^2 .
\]
Thus
\[
        \sup_{\alpha>0}\norm{x^\dag-x_{\alpha,k}}{}
        =
        \norm{e}{}
        =
        \norm{e}{0:k},
\]
which proves inequality (59) at $\nu=0$, in fact with equality.

Now let $\nu=k$.  Then $N_1=1$ and, again by the identity $X_{k:k}=X_k$
from equation (25),
\[
        \norm{e}{k:k}=\norm{e}{k} .
\]
From \eqref{e:erratum:E10:spectral},
\begin{align}
\label{e:erratum:E10:nuk}
        \alpha^{-2k}\norm{S_\alpha^k e}{}^2
        =
        \int_{[0,\infty)}
        (\alpha+\lambda)^{-2k}
        d\mu_e(\lambda).
\end{align}
Since $T$ is injective, $E_{\{0\}}=0$, hence $\mu_e(\{0\})=0$.  For
$\lambda>0$ the function $(\alpha+\lambda)^{-2k}$ increases to
$\lambda^{-2k}$ as $\alpha\searrow0$.  Therefore the monotone convergence
theorem, with the value $+\infty$ allowed, gives
\[
        \lim_{\alpha\searrow0}
        \alpha^{-2k}\norm{S_\alpha^k e}{}^2
        =
        \int_{[0,\infty)} \lambda^{-2k}\,d\mu_e(\lambda)
        =
        \norm{e}{k}^2 .
\]
Consequently,
\[
        \sup_{\alpha>0}
        \alpha^{-k}\norm{x^\dag-x_{\alpha,k}}{}
        =
        \lim_{\alpha\searrow0}
        \alpha^{-k}\norm{x^\dag-x_{\alpha,k}}{}
        =
        \norm{e}{k}
        =
        \norm{e}{k:k} .
\]
This proves inequality (59) at $\nu=k$ and also justifies the final
sentence of Proposition 3.2.
